\chapter{Literature Review}
\label{chap:two}

\section{Introduction}

test

\section{Research gaps and limitations}


\section{unsaturated soil testing techniques}
Filter paper;
HEAV ceremic disc;
pressure plate test;

degree of saturation; gravity moisture content. DEFINITION


\section{Effects of biochar on soil hydraulic properties}
\label{sec:LR-biochar-hydraulic-effects}
The growing demand for more sustainable infrastructure is driving the development of low-carbon, nature-based solutions to soil improvement (DeJong et al. 2013, Ng et al. 2022). For this purpose, biochar has attracted increasing attention as an environmentally friendly amendment because it can reduce soil bulk density, resist microbial degradation, regulate water storage and enhance carbon sequestration (Steinbeiss et al. 2009, Wong et al. 2019, Hussain et al. 2020b, Garg et al. 2020, 2021). This need becomes particularly important as heatwaves and atmospheric aridity intensify evaporative demand and soil moisture depletion (Ng and Pang 2000, Qing et al. 2022, Ng and Peprah-Manu 2023). However, most previous studies on biochar soil improvement have focused on agricultural or vegetated soils, in which biochar is typically used in loosely compacted and frequently irrigated scenarios (Hussain et al. 2020b, 2020a). It remains uncertain whether these hydraulic benefits of biochar can be sustained in geotechnical earthworks (Wong et al. 2018, Hussain et al. 2020b), such as embankments or engineered slopes, where soils may be subjected to much higher compaction and stronger climatic drying (Razouki and Salem 2015, Hussain and Ravi 2022, Walker et al. 2023). Therefore, understanding the coupled effects of biochar and compaction state on soil hydraulic properties is essential to assess its feasibility for geotechnical applications.
The hydraulic response of earthworks to atmospheric drying depends on both the soil-water retention curve (SWRC) and unsaturated hydraulic conductivity (HC) (Ng and Pang 2000, Ng et al. 2023, Hou et al. 2025). This is because SWRC describes how much water remains in the soil as matric suction increases, while HC determines how readily water can move through the unsaturated pore network (Fredlund et al. 1994, Ng and Leung 2012). As drying proceeds, drainage of larger pores reduces the continuity of liquid pathways and limits water transport towards the evaporating surface (An et al. 2018, Hussary et al. 2022). An increase in water retention therefore does not necessarily mean that the rate of moisture loss will be reduced, because drying also depends on the conductivity of the remaining flow paths (Shokri et al. 2008). The addition of biochar further complicates soil hydraulic responses. Its intrapores and surface functional groups can retain water directly, while its low density and angular morphology can alter soil particle packing, pore-size distribution (PSD) and wettability (Suliman et al. 2017, Baiamonte et al. 2019, Mao et al. 2019). Consequently, biochar has been widely demonstrated to enhance soil water retention, whereas its effect on unsaturated HC remains inconsistent. For example, Bagarello et al. (2025) and Villagra-Mendoza and Horn (2018) reported increased unsaturated HC with biochar addition, attributed to a greater proportion and improved connectivity of smaller conducting pores. In contrast, a relevant study by Hussain and Ravi (2021) examined silty sand and pure sand containing 5% and 10% biochar, compacted to 90% of maximum dry density. Biochar enhances water retention capacity but reduces unsaturated HC by approximately one to two orders of magnitude, attributed to pore-size redistribution and greater flow-path tortuosity (Hussain and Ravi 2021). This established a storage-transmission trade-off, but all specimens were examined at one compaction state. In addition, to the best of the authors’ knowledge, relatively few studies have investigated how compaction states, particularly under highly compacted conditions, influence the soil-biochar microstructure and hydraulic behaviours.
Compaction is a key control because it defines the pore space available for biochar accommodation (Ghavanloughajar et al. 2020, Huang et al. 2021). Increasing dry density mainly suppresses larger inter-aggregate pores, modifies drainage throats and alters the air-entry and desaturation response of compacted soils (Su et al. 2022, Ng and Peprah-Manu 2023, Yuan et al. 2024, Hou et al. 2025). Although water retention and permeability have been reported for biochar-amended soils (Wong et al. 2018), the evidence for compaction-dependent SWRC and unsaturated HC is rarely studied (Hussain et al. 2020b, Garg et al. 2021). In particular, by controlling compaction energy, Dong et al. (2024) found that hydrochar increased water retention by filling larger soil pores, while HC could either increase through additional conducting pathways or decrease through pore clogging, depending on dosage. Although hydrochar differs structurally from pyrolysed biochar, these results highlight an important distinction: water retention is governed mainly by redistribution of water-holding pore volume, whereas HC depends more strongly on the continuity, constriction and tortuosity of the active pore-throat network. This distinction is likely to become more important as compaction changes how biochar interacts with the soil fabric (Kameyama et al. 2012). Under loose compaction, biochar may subdivide inter-aggregate voids and create interfacial pores, whereas under dense compaction, the limited free volume may favour throat constriction at aggregate contacts (Baiamonte et al. 2019). The dominant mechanism may also vary with biochar dosage and matric suction as progressively smaller pore domains control water retention and flow during drying. However, the combined effects of biochar dosage and initial compaction state on pore restructuring and the resulting SWRC-HC response during drying remain poorly understood (Hussain et al. 2020b, Dong et al. 2024). This limits mechanistic interpretation of moisture retention and redistribution across loosely and densely compacted earthwork.

\section{Mitigation of microplastics penetration in soil}
The presence of microplastics (MPs) in aquatic systems (Paredes et al., 2019; Yang et al., 2023; Koelmans et al., 2019) has emerged as a global environmental threat (Hale et al., 2020), impacting ecological systems (Pinto da Costa et al., 2019) and human health (Vethaak and Legler, 2021). MP particles, with sizes from \qty{1}{\micro\meter} to \qty{5}{\milli\meter} (Kershaw and Rochman, 2015), are discharged from sources like wastewater (Mason et al., 2016), sewage sludge (Zhang and Chen, 2020) and tyre wear (Tamis et al., 2021), which pollutes soil via surface runoff and rainfall infiltration (Wang et al., 2020). Climate change further accelerated the spread of MP contaminants. For example, extreme weather events such as floods and droughts exacerbate soil surface cracking, thereby intensifying the migration of MP particles into deeper soil layers and water systems (O’Connor et al., 2019). To make the situation even worse, the material and structural characteristics of MPs enable them to serve as vectors for toxins such as pesticides and heavy metals (Kirstein et al., 2016; Koelmans et al., 2016; Prata et al., 2020). After penetrating surface soil (Besseling et al., 2017; Ding et al., 2020) and entering groundwater (Sajjad et al., 2022), MPs imperil the security of the food chain (Hermsen et al., 2017; Nam et al., 2024) and ecological system (Bakir et al., 2014), and are present in biological tissues (Di Fiore et al., 2024) even in the human organs (Zhang et al., 2024). Therefore, it is crucial to understand the mobilisation of MP particles in porous media (e.g. the soil environment), particularly under cyclic hydraulic conditions (i.e., wetting and drying cycles), to prevent their contamination in the deep soil layers and underground water ecosystems (Hsieh et al., 2022; Tong et al., 2020). \par
To address these environmental concerns, it is critical to utilise environmentally friendly and cost-effective materials to prevent the spread of MPs (Zhang et al., 2021b) in agricultural soil, sewage treatment plants, wetlands, and highway corridors (Wang et al., 2024; Fan et al., 2024). Presently, industrial efforts to mitigate the penetration of microplastics (MPs) primarily centre on barrier strategies (Qian et al., 2004), such as compacted clay liners (CCLs), geomembranes, geotextiles and geosynthetic clay liners (GCLs), aimed at offering low permeability against leachates (Koerner, 2005). However, ensuring cost-effective and quality installation poses challenges, while sharp stones frequently puncture these linear structures, leading to unsatisfactory durability and efficacy uncertainty. \par
Biochar, produced from diverse biomass sources, including arboricultural arisings (Bolognesi et al., 2021; Ronsse et al., 2013), has been recognised as an effective, eco-friendly material for wastewater treatment (Gopinath et al., 2021; Berslin et al., 2022; Giorcelli et al., 2019). Biochar can physically retain MP particles by adhering, trapping, and entangling MP spheres (Wang et al., 2020). This capability is attributed to biochar's porous structure (Brewer et al., 2014; Somerville and Jahanshahi, 2015) and large specific surface area (Mohan et al., 2014; Quicker and Weber, 2016; Rajapaksha et al., 2016). Moreover, biochar can actively adsorb MP particles (Kumar et al., 2023), facilitated by reducing electrostatic repulsion or forming chemical bonds between biochar particles and nano- or micro-scaled MPs (Siipola et al., 2020; Wang et al., 2021). Previous studies have demonstrated that MP’s properties and environmental factors can influence their mobilisation in fully saturated conditions (Zhou et al., 2020; Ren et al., 2021). For instance, MPs with smaller sizes, spherical shapes and lower densities exhibit decreased gravitational sedimentation, thus showing higher mobility within the soil (Dong et al., 2021; Dong et al., 2018; Treumann et al., 2014). On the other hand, low ionic strength, significant flow rate, frequent wetting-drying cycles and high porosity can also promote MPs' penetration (Jiang et al., 2021; O’Connor et al., 2019; Wu et al., 2020). \par
However, in an unsaturated soil environment, the migration of MP particles is much more complicated (Liu et al., 2013; Tong et al., 2020). The wetting-drying cycles, acting as external hydraulic loading, can induce the soil’s uneven volumetric strain and crack propagation (Cheng et al., 2020; Tang et al., 2020), which accelerates the vertical migration of MPs into the soil and increases their depth of penetration (O’Connor et al., 2019; Gao et al., 2021). As such, biochar has been demonstrated to enhance the soil's water retention ability, thereby reducing soil desiccation cracks (Lu et al., 2021). This suggests that biochar can potentially prevent MPs from penetrating via soil cracks induced by wetting-drying cycles, not only by actively adsorbing MPs. However, research on unsaturated media has been limited. Even fewer studies have investigated the MP-retention performance of biochar-amended soil under wetting-drying cycles to demonstrate the reliability of biochar amendment in mitigating the effects of hydraulic factors. Regarding the combination of biochar and soil, there has been no uniform format, either by direct mixing (Sun et al., 2018; Wang et al., 2022) or layering (Wang et al., 2020; Ahmad et al., 2023). Few studies compared these two formats to optimise a better form of biochar and soil. Therefore, this application of biochar and the underlying mechanisms need to be further investigated to support this potential industry application.\par
To demonstrate the application of biochar in mitigating the negative effect of the wetting-drying cycle and to optimise the mixing format of biochar and soil, this study performed a series of column tests considering different biochar contents, initial void ratios and biochar-soil mixing formats. Utilising breakthrough curves, water retention curves and MPs distribution patterns, the MPs' retention capability and crack resistance under wetting-drying cycles were analysed. Scanning electron microscopy (SEM) and Fourier-transform infrared spectroscopy (FTIR) methods were employed to investigate the morphology of particles and analyse the adsorption mechanisms. Finally, recommendations for potential practical applications, such as soil-biochar mixing format, were proposed. \par

\section{Summary}